\section{Reversible critical weights on the original collision record}
\label{sec:v65-reversible-critical}
The positive source retained near a clearance caustic is sensitive to
critical coefficients, not just to the size of a roof neighborhood.
We first identify the coefficient of a normal-to-normal word with a
reversible linearization. The identity holds at every regular word,
without a lower bound on its interior incidences. At a simple
competing-hit boundary it has a physical-side interpretation, which
must be distinguished from the derivative of the singular billiard map.

Determinant identities relating action Hessians and monodromy are
classical; the discrete and reversible theory is developed in
\cite{BT65}. The purpose here is the exact source normalization, the
half-word labels, and the square-root weight needed by the roof density.
No periodic-orbit equidistribution theorem is inferred from the identity.

\subsection{Normal endpoints and reversible closure}
Fix a regular word of $m\ge1$ flights contributing to the original
section source. Its initial and terminal section tests and its exact
half-word labels are retained. Use canonical endpoint
arclengths $u,v$, with tangential momenta $p_0,p_m$, and let $F(u,v)$
be its reduced action. Suppose $z$ is a critical point of $F$.
The first variation gives $p_0=p_m=0$. Write
\begin{equation}\label{eq:v65-hessian-data}
 H_z=\nabla^2F(z)=
 \begin{pmatrix}h_{00}&h_{01}\\h_{01}&h_{11}\end{pmatrix},\qquad
 M_z=D T_R^m(z)=\begin{pmatrix}a&b\\c&d\end{pmatrix}.
\end{equation}
The matrix entry $c$ in this display is not the section mass; the
latter is denoted throughout this section by $c_*=91/(10000\pi)$.
Matrices use arclength and momentum, not angle and momentum with an
unrecorded scale factor.

\begin{lemma}[The half-word and its reflected closure]
\label{lem:v65-reversible-matrix}
The point $z$ is fixed by $T_R^{2m}$ on the collision quotient, and
\begin{equation}\label{eq:v65-reversible-matrix}
 P_z:=D T_R^{2m}(z)=S M_z^{-1}S M_z,
       \qquad S=\operatorname{diag}(1,-1).
\end{equation}
Moreover
\begin{align}\label{eq:v65-generating-matrix}
 b=-h_{01}^{-1},\quad a=-h_{00}/h_{01},\quad
 d=-h_{11}/h_{01},\quad c=-\det H_z/h_{01},
 \qquad bc=\det H_z/h_{01}^2>0,
\end{align}
and
\begin{equation}\label{eq:v65-monodromy-determinant}
 P_z=\begin{pmatrix}ad+bc&2bd\\2ac&ad+bc\end{pmatrix},\qquad
 \det(I-P_z)=-4bc.
\end{equation}
The period need not be primitive. There is no division by $m$ or
$2m$ in these assertions.
\end{lemma}
\begin{proof}
The collision reversal is $\iota(u,p)=(u,-p)$ and satisfies
$\iota T_R\iota=T_R^{-1}$ on regular orbits. Both $z$ and
$z_m=T_R^m z$ are fixed by $\iota$. Thus
$T_R^m z_m=\iota T_R^{-m}\iota z_m=z$; geometrically the second
half retraces the first. Differentiating the reversal identity at
$z_m$ gives $D T_R^m(z_m)=S M_z^{-1}S$, proving
\eqref{eq:v65-reversible-matrix}.

Differentiate $p_0=-F_u$ and $p_m=F_v$ and solve the first equation
for $dv$ in terms of $du,dp_0$. This gives every entry in
\eqref{eq:v65-generating-matrix}. In particular $\det M_z=1$.
The endpoint injectivity and nonzero mixed derivative are those of
Lemma~\ref{lem:v44-endpoint-curvature} and
\eqref{eq:v61-cross-and-product}. Positivity of $H_z$ gives $bc>0$.
Multiplication, using $ad-bc=1$, gives
\eqref{eq:v65-monodromy-determinant}.
\end{proof}

The closure is used only to calculate a weight. The displacement,
return index and roof of the original source are still those of the
first $m$ flights. In particular the section is not assumed invariant
under reversal. Neither its occupation count nor its arithmetic
residue is read from the doubled orbit. Two starting points on one
closed orbit are not identified if they label distinct original
source states.

\begin{theorem}[The intrinsic coefficient as a reversible weight]
\label{thm:v65-reversible-weight}
Let $J_z$ be the full-word Morse coefficient of the common density
$\nu/c_*$, before any active one-sided restriction. At a strict
section point this is \eqref{eq:v45-intrinsic-jump} for the original
source; boundary fractions will be imposed in the next two sections.
The eigenvalues of $P_z$ are $\lambda_z,\lambda_z^{-1}$ with
$\lambda_z>1$, and
\begin{equation}\label{eq:v65-reversible-weight}
 J_z=\frac{1}{R c_*\sqrt{|\det(I-P_z)|}}
     =\frac{1}{R c_*}
        \frac{\lambda_z^{-1/2}}{1-\lambda_z^{-1}}.
\end{equation}
There are fixed $C<\infty$ and $q_0=47/53<1$ such that
\begin{equation}\label{eq:v65-single-weight-decay}
 J_z\le Cq_0^{m-2}\quad(m\ge2),
\end{equation}
uniformly in $R$ and in all positive interior incidences.
The square-root determinant in \eqref{eq:v65-reversible-weight}
cannot be replaced by the ordinary inverse determinant.
\end{theorem}
\begin{proof}
The original endpoint density at $z$ is
$w_z=|h_{01}|/(4\pi R c_*)$. Hence the two-dimensional Morse
coefficient is
\[
 J_z=\frac{2\pi w_z}{\sqrt{\det H_z}}
     =\frac{1}{2R c_*\sqrt{bc}}.
\]
Equation~\eqref{eq:v65-monodromy-determinant} gives the first formula.
Since $\operatorname{tr}P_z=2+4bc>2$ and $\det P_z=1$, its
eigenvalues are positive reciprocal numbers, one greater than one.
Now $|\det(I-P_z)|=(\lambda_z-1)^2/\lambda_z$, which proves the
second formula and fixes its square root.

For the count bound use the internal matrix $A$ in
\eqref{eq:v45-contact-matrix}. The endpoint Schur complement gives,
for $m\ge2$ and normal endpoints,
\[
 |h_{01}|=t_0t_{m-1}|(A^{-1})_{1,m-1}|.
\]
This includes the one-entry matrix when $m=2$. The Neumann estimate
\eqref{eq:v45-inverse-contact}, valid for every positive incidence,
bounds this by $Cq_0^{m-2}$. Also
$H_z\ge R^{-1}I$ in arclength coordinates by the endpoint second
variation. Thus $\sqrt{\det H_z}\ge R^{-1}$ and the formula for
$J_z$ proves \eqref{eq:v65-single-weight-decay}. This proof does
not bound a sum over words by its largest summand.
\end{proof}

\subsection{A normal-section crossing measure}
There are constants $0<\lambda_*\le\Lambda_*<\infty$, independent
of the word and of all its positive incidences, such that at a normal
critical point
\begin{equation}\label{eq:v65-normal-hessian-bounds}
 \lambda_* I\le H_z\le\Lambda_* I.
\end{equation}
One may take $\lambda_*=1/R_+$ and
$\Lambda_*=2(1/R_-+1/\ell_0)$. The lower bound was just proved.
The internal Schur complement is bounded above by the endpoint
principal block. Its entries are bounded by endpoint curvature and
$1/\ell_0$; the stated choice also includes $m=1$.

\begin{proposition}[Crossing weights and fixed-count traces]
\label{prop:v65-normal-crossing}
Set $W_z=|\partial_u p_m(u,0)|_z^{-1}=|c|^{-1}$. Then
\begin{equation}\label{eq:v65-crossing-comparison}
 J_z=\frac{\sqrt{\det H_z}}{2R c_*}W_z,
 \qquad
 \frac{\lambda_*}{2R_+c_*}W_z\le J_z
       \le\frac{\Lambda_*}{2R_-c_*}W_z.
\end{equation}
For a regular critical point whose physical and section inequalities
are all strict, take an isolating source neighborhood $U_z$.
For any continuous roof test $\psi$,
\begin{equation}\label{eq:v65-normal-trace-limit}
 \lim_{d\downarrow0}\frac{1}{4d^2}
 \int_{U_z}\one_{\{|p_0|<d,\ |p_m|<d\}}
             \psi(L_{m,R})\,d\nu_R^*
       =\frac{W_z}{4\pi R c_*}\psi(t_z),\qquad t_z=F(z).
\end{equation}
The collision count in this limit is fixed.
\end{proposition}
\begin{proof}
The ratio $c/b=\det H_z$ in
\eqref{eq:v65-generating-matrix} proves the equality in
\eqref{eq:v65-crossing-comparison}; then apply
\eqref{eq:v65-normal-hessian-bounds}. The map
$(u,p_0)\mapsto(p_0,p_m)$ has absolute determinant $|c|$ at $z$.
It is a local diffeomorphism. Its inverse sends the shrinking square
$(-d,d)^2$ to $U_z$ for all sufficiently small $d$. The source
density is $1/(4\pi R c_*)$ in $(u,p_0)$. Change variables,
divide by the square's area, and use continuity of the integrand.
\end{proof}

For a finite family with the same original labels define
\begin{equation}\label{eq:v65-reversible-atomic-measure}
 \mathfrak T_{n,k,m,R}
   =\sum_z\frac{\lambda_z^{-1/2}}{1-\lambda_z^{-1}}\delta_{t_z}.
 \qquad
 \sum_zJ_z\delta_{t_z}=(R c_*)^{-1}\mathfrak T_{n,k,m,R}.
\end{equation}
This is a positive atomic measure, not a probability or a spectral
trace of an operator already constructed in the paper. Equal roof
values are added. The pointwise identity gives equality of its
weighted sums against every bounded test, so it does not require
separation of critical values.

\subsection{Physical-side values at a clearance seam}
\begin{proposition}[The boundary interpretation]
\label{prop:v65-boundary-monodromy}
Suppose a selected collision word has a physical-side limit $z$ at
a competing-hit seam, with every selected contact incidence positive
and both endpoint momenta zero. Its simple contact equations define
local smooth $F$, $H_z$, $M_z$, and $w_z$ as limits of that branch.
Define $P_z=S M_z^{-1}S M_z$ algebraically. Equations
\eqref{eq:v65-monodromy-determinant},
\eqref{eq:v65-reversible-weight},
\eqref{eq:v65-single-weight-decay} and
\eqref{eq:v65-crossing-comparison} remain valid for these branch data.
They do not assert that $D T_R^{2m}$ exists at the singular state.
\end{proposition}
\begin{proof}
Positive selected incidence makes each contact root simple. The
implicit contact graph and its endpoint map therefore extend smoothly
locally, whether or not the strict competing-hit inequality survives
on both sides. The Hessian, mixed derivative and symplectic generating
identities persist by continuity from the physical side. The positive
normal Hessian and the nonzero mixed derivative persist at this finite
word, by the explicit contact determinant. Apply the same algebra.
The extension is used to compute derivatives only; no source is placed
on a nonphysical continuation. At the seam the actual collision map
can change its selected first hit, and its full derivative is not used.
\end{proof}

\paragraph{Relation to the ordered height problem.}
The exact formula identifies a concrete positive weight to be summed
in the original labels. It does not prove the required long-count
concentration estimate for that sum. In particular the limit
$d\downarrow0$ in \eqref{eq:v65-normal-trace-limit} cannot be
interchanged with the collision limit in
Lemma~\ref{lem:v45-two-strip-local}. That interchange would require
a uniform transverse-trace estimate not supplied by the fixed-width
local theorem. Likewise individual exponential decay in
\eqref{eq:v65-single-weight-decay} is not a bound on an exponentially
large critical cluster. The next section uses the exact weights to
calculate, rather than discard, a simple roof-critical clearance source.
